\begin{trapbox}

	\begin{description}[style=nextline, leftmargin=2.8em, labelsep=0.5em, itemsep=0.35em]
		\item[\textbf{\textcolor{CTrap}{易错点一（漏乘系数）}}]
			计算锥体体积时，误使用$V=Sh$，漏写因子$\frac{1}{3}$。
		\item[\textbf{\textcolor{CTrap}{正确理解（系数来历）}}]
			锥体体积是等底等高柱体体积的三分之一，公式为$V=\frac{1}{3}Sh$，系数$\frac{1}{3}$不可少。
		\item[\textbf{\textcolor{CTrap}{错因分析（混淆形状）}}]
			将锥体与柱体体积公式混淆，或凭直觉以为锥体是柱体的一半（实际为三分之一）。
	\end{description}

	\medskip
	\begin{description}[style=nextline, leftmargin=2.8em, labelsep=0.5em, itemsep=0.35em]
		\item[\textbf{\textcolor{CTrap}{易错点二（底面积算错）}}]
			计算底面面积时，用了错误的数据或公式，例如把圆锥的直径当作半径，或棱锥底面多边形面积计错。
		\item[\textbf{\textcolor{CTrap}{正确理解（核对公式）}}]
			根据底面形状使用对应面积公式：圆用$\pi r^2$，正三角形用$\frac{\sqrt{3}}{4}a^2$，正方形用$a^2$，梯形用$\frac{(a+b)h}{2}$等，并仔细核对已知条件。
		\item[\textbf{\textcolor{CTrap}{错因分析（粗心失误）}}]
			对面积公式记忆不牢，或审题时看错已知量（如半径与直径混淆）。
	\end{description}

	\medskip
	\begin{description}[style=nextline, leftmargin=2.8em, labelsep=0.5em, itemsep=0.35em]
		\item[\textbf{\textcolor{CTrap}{易错点三（高取错）}}]
			将棱锥的斜高、侧棱长或圆锥的母线当作高代入公式。
		\item[\textbf{\textcolor{CTrap}{正确理解（垂直高）}}]
			高是从顶点向底面作垂线的垂直距离。在直角锥中，高可通过勾股定理由侧棱和底面外接圆半径等求得。
		\item[\textbf{\textcolor{CTrap}{错因分析（概念混淆）}}]
			不熟悉几何体的高、斜高、母线等概念的区别，或解题时未从条件中正确求出高。
	\end{description}

\end{trapbox}
